% ============================================================
% CHƯƠNG: PHÂN TÍCH YÊU CẦU HỆ THỐNG — E-MAPP
% Dịch vụ công số tỉnh Thanh Hóa
% ============================================================

\section{Phân tích yêu cầu hệ thống}

Chương này trình bày kết quả phân tích yêu cầu đối với hệ thống E-Mapp —
cổng dịch vụ công số tỉnh Thanh Hóa, bao gồm các yêu cầu chức năng,
yêu cầu phi chức năng và định nghĩa các tác nhân tham gia hệ thống.

% ─────────────────────────────────────────────────────────────
\subsection{Yêu cầu chức năng (Functional Requirements)}
% ─────────────────────────────────────────────────────────────

Yêu cầu chức năng mô tả các nghiệp vụ mà hệ thống phải thực hiện được.
Toàn bộ yêu cầu được phân nhóm theo các module chức năng chính.

% ── FR-1 ─────────────────────────────────────────────────────
\subsubsection{FR-1: Quản lý xác thực và tài khoản người dùng}

\begin{table}[H]
\centering
\caption{Yêu cầu chức năng — Xác thực người dùng (FR-1)}
\label{tab:fr-auth}
\begin{tabularx}{\textwidth}{|c|l|X|c|}
\hline
\textbf{Mã} & \textbf{Tên chức năng} & \textbf{Mô tả} & \textbf{Ưu tiên} \\
\hline
FR-1.1 & Đăng ký tài khoản &
  Công dân đăng ký bằng số CCCD (12 chữ số), họ tên, ngày sinh, số điện thoại,
  email và mật khẩu đáp ứng chính sách bảo mật (≥8 ký tự, có chữ hoa, số,
  ký tự đặc biệt). Xác thực danh tính qua mã OTP 6 chữ số. &
  Cao \\
\hline
FR-1.2 & Xác thực qua VNeID &
  Hỗ trợ xác thực định danh điện tử qua hệ thống VNeID quốc gia.
  Sau khi xác thực thành công, tài khoản được đánh dấu \texttt{isVNeIDVerified = true}
  và lưu mã định danh VNeID. &
  Cao \\
\hline
FR-1.3 & Đăng nhập &
  Người dùng đăng nhập bằng số CCCD và mật khẩu. Hệ thống trả về
  JWT token lưu phía client (\texttt{localStorage}), có thời hạn hiệu lực
  cấu hình được. &
  Cao \\
\hline
FR-1.4 & Đăng xuất &
  Vô hiệu hóa phiên làm việc, xóa JWT token khỏi client. &
  Cao \\
\hline
FR-1.5 & Quên mật khẩu &
  Công dân có thể đặt lại mật khẩu qua email hoặc số điện thoại đã đăng ký. &
  Trung bình \\
\hline
FR-1.6 & Xem hồ sơ cá nhân &
  Hiển thị thông tin tài khoản: họ tên, CCCD, ngày sinh, SĐT, email,
  trạng thái xác thực VNeID. &
  Cao \\
\hline
FR-1.7 & Cập nhật hồ sơ cá nhân &
  Cho phép sửa họ tên, SĐT, email, ngày sinh. Dữ liệu được validate
  trước khi lưu. &
  Trung bình \\
\hline
\end{tabularx}
\end{table}

% ── FR-2 ─────────────────────────────────────────────────────
\subsubsection{FR-2: Dịch vụ công và nộp hồ sơ trực tuyến}

\begin{table}[H]
\centering
\caption{Yêu cầu chức năng — Dịch vụ công \& Hồ sơ (FR-2)}
\label{tab:fr-services}
\begin{tabularx}{\textwidth}{|c|l|X|c|}
\hline
\textbf{Mã} & \textbf{Tên chức năng} & \textbf{Mô tả} & \textbf{Ưu tiên} \\
\hline
FR-2.1 & Xem danh mục thủ tục &
  Hiển thị danh mục thủ tục hành chính phân theo cấp (xã/phường, quận/huyện,
  tỉnh) và theo lĩnh vực. Hỗ trợ lọc đa tiêu chí. &
  Cao \\
\hline
FR-2.2 & Tìm kiếm thủ tục &
  Tìm kiếm thủ tục theo từ khóa, danh mục, cấp hành chính. &
  Cao \\
\hline
FR-2.3 & Xem chi tiết thủ tục &
  Hiển thị đầy đủ thông tin: yêu cầu hồ sơ, biểu mẫu, thời hạn xử lý,
  lệ phí, cơ quan thụ lý. &
  Cao \\
\hline
FR-2.4 & Nộp hồ sơ trực tuyến &
  Công dân điền biểu mẫu và nộp hồ sơ điện tử. Hỗ trợ lưu bản nháp
  (\textit{draft}) trước khi nộp chính thức. &
  Cao \\
\hline
FR-2.5 & Tải lên tài liệu &
  Cho phép đính kèm file PDF, ảnh (JPG/PNG/GIF) tối đa 16~MB/file.
  Hỗ trợ nhiều file trong một hồ sơ. &
  Cao \\
\hline
FR-2.6 & Ký điện tử xác nhận &
  Hỗ trợ ký điện tử / xác nhận hồ sơ trước khi nộp chính thức. &
  Trung bình \\
\hline
FR-2.7 & Tra cứu trạng thái hồ sơ &
  Công dân theo dõi trạng thái hồ sơ theo luồng: \textit{Bản nháp → Đã nộp
  → Đang xem xét → Đã duyệt / Từ chối / Yêu cầu bổ sung}. &
  Cao \\
\hline
FR-2.8 & Rút hồ sơ &
  Công dân có thể rút lại hồ sơ đã nộp khi chưa được xử lý. &
  Trung bình \\
\hline
FR-2.9 & Bổ sung hồ sơ &
  Khi cán bộ yêu cầu bổ sung, công dân nhận thông báo và nộp
  tài liệu bổ sung vào hồ sơ hiện có. &
  Cao \\
\hline
FR-2.10 & Xử lý hồ sơ (Admin) &
  Cán bộ/admin xem danh sách, duyệt, từ chối (có ghi chú), hoặc
  yêu cầu bổ sung thông tin hồ sơ. &
  Cao \\
\hline
FR-2.11 & AI trích xuất tài liệu &
  Hệ thống AI tự động trích xuất thông tin từ file tài liệu đính kèm
  (OCR + NLP) để hỗ trợ xét duyệt. &
  Trung bình \\
\hline
\end{tabularx}
\end{table}

% ── FR-3 ─────────────────────────────────────────────────────
\subsubsection{FR-3: Đặt lịch hẹn}

\begin{table}[H]
\centering
\caption{Yêu cầu chức năng — Đặt lịch hẹn (FR-3)}
\label{tab:fr-appointment}
\begin{tabularx}{\textwidth}{|c|l|X|c|}
\hline
\textbf{Mã} & \textbf{Tên chức năng} & \textbf{Mô tả} & \textbf{Ưu tiên} \\
\hline
FR-3.1 & Đặt lịch hẹn trực tuyến &
  Công dân chọn cơ quan, loại dịch vụ, ngày và khung giờ còn trống để
  đặt lịch hẹn. Hệ thống kiểm tra slot khả dụng và ngăn đặt lịch trong quá khứ. &
  Cao \\
\hline
FR-3.2 & Xem lịch hẹn dạng Calendar &
  Hiển thị lịch hẹn của công dân trên giao diện lịch (FullCalendar),
  hỗ trợ xem theo tuần và tháng. &
  Cao \\
\hline
FR-3.3 & Xem danh sách lịch hẹn &
  Liệt kê toàn bộ lịch hẹn với trạng thái: \textit{pending / confirmed / cancelled}. &
  Cao \\
\hline
FR-3.4 & Hủy lịch hẹn &
  Cho phép hủy lịch hẹn chưa đến hạn. &
  Trung bình \\
\hline
FR-3.5 & Đặt lịch bằng giọng nói &
  Công dân nói yêu cầu bằng tiếng Việt; AI thực hiện hội thoại
  multi-turn (hỏi cơ quan, dịch vụ, ngày, giờ) rồi tự động tạo lịch hẹn. &
  Cao \\
\hline
FR-3.6 & Xem lịch hẹn theo cơ quan (Admin) &
  Cán bộ xem danh sách lịch hẹn của cơ quan mình theo ngày. &
  Trung bình \\
\hline
\end{tabularx}
\end{table}

% ── FR-4 ─────────────────────────────────────────────────────
\subsubsection{FR-4: Hệ thống hàng chờ}

\begin{table}[H]
\centering
\caption{Yêu cầu chức năng — Hàng chờ (FR-4)}
\label{tab:fr-queue}
\begin{tabularx}{\textwidth}{|c|l|X|c|}
\hline
\textbf{Mã} & \textbf{Tên chức năng} & \textbf{Mô tả} & \textbf{Ưu tiên} \\
\hline
FR-4.1 & Lấy số thứ tự &
  Công dân (đã đăng nhập hoặc khách) lấy số thứ tự tại cơ quan.
  Vé được tạo tự động với số thứ tự tuần tự trong ngày. &
  Cao \\
\hline
FR-4.2 & Xem trạng thái vé &
  Hiển thị vị trí trong hàng, ước tính thời gian chờ,
  trạng thái hiện tại của vé. &
  Cao \\
\hline
FR-4.3 & Xem vé của tôi hôm nay &
  Công dân đã đăng nhập xem toàn bộ vé đã lấy trong ngày. &
  Cao \\
\hline
FR-4.4 & Hủy vé hàng chờ &
  Công dân hủy vé chưa được phục vụ. &
  Trung bình \\
\hline
FR-4.5 & Gọi số tiếp theo (Staff) &
  Nhân viên gọi vé tiếp theo tại quầy.
  Hệ thống broadcast cập nhật realtime qua WebSocket. &
  Cao \\
\hline
FR-4.6 & Cập nhật trạng thái vé (Staff) &
  Nhân viên cập nhật trạng thái: \textit{waiting / called / serving /
  done / absent / cancelled}. &
  Cao \\
\hline
FR-4.7 & Quản lý quầy dịch vụ &
  Admin/Staff tạo và cập nhật thông tin quầy giao dịch của cơ quan. &
  Trung bình \\
\hline
FR-4.8 & Màn hình hiển thị hàng chờ &
  Màn hình TV tại cơ quan hiển thị số đang phục vụ, số đang chờ;
  cập nhật realtime qua WebSocket. &
  Cao \\
\hline
FR-4.9 & Snapshot hàng chờ trên bản đồ &
  Cung cấp API công khai trả về tóm tắt hàng chờ tất cả cơ quan
  để hiển thị trên bản đồ. &
  Trung bình \\
\hline
\end{tabularx}
\end{table}

% ── FR-5 ─────────────────────────────────────────────────────
\subsubsection{FR-5: Bản đồ và định vị}

\begin{table}[H]
\centering
\caption{Yêu cầu chức năng — Bản đồ (FR-5)}
\label{tab:fr-map}
\begin{tabularx}{\textwidth}{|c|l|X|c|}
\hline
\textbf{Mã} & \textbf{Tên chức năng} & \textbf{Mô tả} & \textbf{Ưu tiên} \\
\hline
FR-5.1 & Xem bản đồ cơ quan &
  Hiển thị vị trí các cơ quan hành chính trên nền bản đồ OpenStreetMap. &
  Cao \\
\hline
FR-5.2 & Tìm kiếm địa điểm &
  Tra cứu địa điểm với gợi ý tự động (autocomplete) sử dụng
  Nominatim API. &
  Cao \\
\hline
FR-5.3 & Geocode / Reverse geocode &
  Chuyển đổi địa chỉ sang tọa độ và ngược lại. &
  Cao \\
\hline
FR-5.4 & Tìm đường đến cơ quan &
  Tính toán và hiển thị lộ trình từ vị trí người dùng đến cơ quan
  sử dụng OSRM (miễn phí, không cần API key). &
  Cao \\
\hline
FR-5.5 & Tìm cơ quan gần nhất &
  Lọc và hiển thị các cơ quan/dịch vụ trong bán kính chỉ định
  quanh vị trí hiện tại của người dùng. &
  Trung bình \\
\hline
FR-5.6 & Xem hàng chờ trên bản đồ &
  Overlay trạng thái hàng chờ realtime lên bản đồ để người dùng
  chọn cơ quan ít đông nhất. &
  Trung bình \\
\hline
\end{tabularx}
\end{table}

% ── FR-6 ─────────────────────────────────────────────────────
\subsubsection{FR-6: AI Chatbot, RAG và Voice}

\begin{table}[H]
\centering
\caption{Yêu cầu chức năng — AI \& Voice (FR-6)}
\label{tab:fr-ai}
\begin{tabularx}{\textwidth}{|c|l|X|c|}
\hline
\textbf{Mã} & \textbf{Tên chức năng} & \textbf{Mô tả} & \textbf{Ưu tiên} \\
\hline
FR-6.1 & Chatbot hỏi đáp thủ tục &
  Hỗ trợ hội thoại tiếng Việt về các thủ tục hành chính.
  Chatbot trả lời dựa trên cơ sở tri thức RAG. &
  Cao \\
\hline
FR-6.2 & Tìm kiếm RAG &
  Hệ thống Retrieval-Augmented Generation sử dụng embedding tiếng Việt
  để tìm tài liệu liên quan và sinh câu trả lời có ngữ cảnh chính xác. &
  Cao \\
\hline
FR-6.3 & Gợi ý thủ tục phù hợp &
  AI phân tích nhu cầu của công dân và đề xuất các thủ tục hành chính
  phù hợp nhất. &
  Trung bình \\
\hline
FR-6.4 & Lịch sử cuộc trò chuyện &
  Lưu và hiển thị lịch sử chat theo phiên (session), hỗ trợ tiếp tục
  hội thoại từ lần trước. &
  Trung bình \\
\hline
FR-6.5 & Nhận diện giọng nói (STT) &
  Chuyển đổi giọng nói tiếng Việt sang văn bản sử dụng Google
  Cloud Speech-to-Text. &
  Cao \\
\hline
FR-6.6 & Tổng hợp giọng nói (TTS) &
  Chuyển đổi văn bản phản hồi sang âm thanh tiếng Việt tự nhiên
  sử dụng Google Cloud Text-to-Speech. &
  Cao \\
\hline
FR-6.7 & Hội thoại đặt lịch (Voice Dialog) &
  Pipeline Audio → STT → NLU (Gemini) → phản hồi → TTS.
  Hội thoại multi-turn tự động thu thập thông tin đặt lịch. &
  Cao \\
\hline
FR-6.8 & Cấu hình chatbot (Admin) &
  Admin tùy chỉnh system prompt, hành vi chatbot và reset về mặc định. &
  Thấp \\
\hline
\end{tabularx}
\end{table}

% ── FR-7 ─────────────────────────────────────────────────────
\subsubsection{FR-7: Quản trị hệ thống}

\begin{table}[H]
\centering
\caption{Yêu cầu chức năng — Quản trị (FR-7)}
\label{tab:fr-admin}
\begin{tabularx}{\textwidth}{|c|l|X|c|}
\hline
\textbf{Mã} & \textbf{Tên chức năng} & \textbf{Mô tả} & \textbf{Ưu tiên} \\
\hline
FR-7.1 & Dashboard thống kê &
  Bảng điều khiển tổng hợp số liệu: số hồ sơ theo trạng thái,
  số lịch hẹn, lưu lượng hàng chờ, biểu đồ xu hướng. &
  Cao \\
\hline
FR-7.2 & Quản lý tài khoản người dùng &
  Xem danh sách, chi tiết tài khoản công dân; quản lý phân quyền
  (\textit{citizen / staff / admin}). &
  Cao \\
\hline
FR-7.3 & Quản lý địa điểm / cơ quan &
  CRUD danh sách địa điểm, thông tin liên hệ, tọa độ, giờ làm việc
  của các cơ quan hành chính. &
  Cao \\
\hline
FR-7.4 & Quản lý danh mục thủ tục &
  Thêm, sửa, xóa danh mục và chi tiết thủ tục hành chính. &
  Cao \\
\hline
FR-7.5 & Xem log hệ thống &
  Theo dõi log hoạt động, lỗi hệ thống theo module. &
  Trung bình \\
\hline
\end{tabularx}
\end{table}

% ── FR-8 ─────────────────────────────────────────────────────
\subsubsection{FR-8: Đánh giá và thông báo}

\begin{table}[H]
\centering
\caption{Yêu cầu chức năng — Đánh giá \& Thông báo (FR-8)}
\label{tab:fr-eval}
\begin{tabularx}{\textwidth}{|c|l|X|c|}
\hline
\textbf{Mã} & \textbf{Tên chức năng} & \textbf{Mô tả} & \textbf{Ưu tiên} \\
\hline
FR-8.1 & Đánh giá chất lượng dịch vụ &
  Công dân đánh giá mức độ hài lòng (thang 1–5 sao) và gửi
  nhận xét văn bản sau khi hoàn thành dịch vụ. &
  Trung bình \\
\hline
FR-8.2 & Xem trung tâm thông báo &
  Hiển thị danh sách thông báo: cập nhật trạng thái hồ sơ,
  nhắc nhở lịch hẹn, thông báo hệ thống. &
  Cao \\
\hline
FR-8.3 & Đánh dấu thông báo đã đọc &
  Cho phép đánh dấu từng thông báo hoặc tất cả là đã đọc. &
  Trung bình \\
\hline
\end{tabularx}
\end{table}


% ─────────────────────────────────────────────────────────────
\subsection{Yêu cầu phi chức năng}
% ─────────────────────────────────────────────────────────────

Yêu cầu phi chức năng xác định các ràng buộc về chất lượng mà hệ thống
phải đáp ứng, bao gồm ba nhóm chính: hiệu năng, tính khả dụng và bảo mật.

% ── NFR-1: Hiệu năng ─────────────────────────────────────────
\subsubsection{NFR-1: Hiệu năng (Performance)}

\begin{table}[H]
\centering
\caption{Yêu cầu phi chức năng — Hiệu năng}
\label{tab:nfr-perf}
\begin{tabularx}{\textwidth}{|c|l|X|}
\hline
\textbf{Mã} & \textbf{Chỉ tiêu} & \textbf{Mức yêu cầu} \\
\hline
NFR-1.1 & Thời gian phản hồi API &
  Các API CRUD thông thường phải phản hồi trong vòng \textbf{500ms}
  ở điều kiện tải bình thường (≤ 100 người dùng đồng thời). \\
\hline
NFR-1.2 & Thời gian phản hồi AI Chatbot &
  Câu trả lời từ RAG chatbot phải có trong vòng \textbf{3 giây}
  (bao gồm cả thời gian tìm kiếm embedding và sinh câu trả lời). \\
\hline
NFR-1.3 & Thời gian xử lý STT/TTS &
  Pipeline Speech-to-Text hoàn thành trong \textbf{2 giây} cho
  đoạn âm thanh dưới 15 giây. \\
\hline
NFR-1.4 & Tải trang web &
  Giao diện React (đã build) tải hoàn toàn trong \textbf{2 giây}
  trên kết nối 4G (First Contentful Paint ≤ 1.5s). \\
\hline
NFR-1.5 & WebSocket realtime &
  Cập nhật hàng chờ realtime phải được đẩy đến client trong
  \textbf{200ms} sau khi có sự kiện thay đổi. \\
\hline
NFR-1.6 & Tải lên tài liệu &
  Upload file ≤ 16MB phải hoàn thành trong \textbf{10 giây}
  trên đường truyền 10~Mbps. \\
\hline
NFR-1.7 & Cache AI (CAG) &
  Hệ thống CAG cache phải giảm thời gian trả lời cho câu hỏi
  đã được hỏi trước đó xuống \textbf{≤ 100ms}. \\
\hline
\end{tabularx}
\end{table}

% ── NFR-2: Tính khả dụng ─────────────────────────────────────
\subsubsection{NFR-2: Tính khả dụng (Availability \& Reliability)}

\begin{table}[H]
\centering
\caption{Yêu cầu phi chức năng — Tính khả dụng}
\label{tab:nfr-avail}
\begin{tabularx}{\textwidth}{|c|l|X|}
\hline
\textbf{Mã} & \textbf{Chỉ tiêu} & \textbf{Mức yêu cầu} \\
\hline
NFR-2.1 & Độ sẵn sàng hệ thống &
  Hệ thống phải đạt \textbf{SLA 99,5\%} uptime trong giờ hành chính
  (7:00 – 17:00 các ngày làm việc). \\
\hline
NFR-2.2 & Thời gian phục hồi &
  Thời gian phục hồi sau sự cố (\textit{Recovery Time Objective})
  không quá \textbf{30 phút}. \\
\hline
NFR-2.3 & Hỗ trợ đa thiết bị &
  Giao diện hoạt động đúng trên trình duyệt Chrome, Firefox, Edge
  phiên bản mới nhất; tương thích màn hình desktop và mobile
  (responsive design). \\
\hline
NFR-2.4 & Khả năng mở rộng &
  Kiến trúc backend (Flask + SQLAlchemy) hỗ trợ scale ngang
  (\textit{horizontal scaling}) khi số lượng người dùng tăng.
  Cơ sở dữ liệu PostgreSQL/SQL Server hỗ trợ connection pooling. \\
\hline
NFR-2.5 & Tính liên tục dữ liệu &
  Dữ liệu hồ sơ, lịch hẹn và hàng chờ phải được đồng bộ giữa
  bộ nhớ Redis/in-memory và cơ sở dữ liệu quan hệ. Không mất dữ liệu
  khi khởi động lại server. \\
\hline
NFR-2.6 & Xử lý lỗi graceful &
  Mọi lỗi phía server phải trả về HTTP code phù hợp và thông điệp
  tiếng Việt rõ ràng. Lỗi không được để lộ stack trace ra client. \\
\hline
NFR-2.7 & Ghi log đầy đủ &
  Mọi request và sự kiện quan trọng phải được ghi log có cấu trúc
  (module, level, timestamp) để hỗ trợ giám sát và debug. \\
\hline
\end{tabularx}
\end{table}

% ── NFR-3: Bảo mật ───────────────────────────────────────────
\subsubsection{NFR-3: Bảo mật (Security)}

\begin{table}[H]
\centering
\caption{Yêu cầu phi chức năng — Bảo mật}
\label{tab:nfr-sec}
\begin{tabularx}{\textwidth}{|c|l|X|}
\hline
\textbf{Mã} & \textbf{Chỉ tiêu} & \textbf{Mức yêu cầu} \\
\hline
NFR-3.1 & Xác thực bằng JWT &
  Tất cả API cần bảo vệ phải yêu cầu JWT token hợp lệ trong header.
  Token được ký bằng thuật toán \texttt{HS256} với khóa bí mật cấu hình
  qua biến môi trường. \\
\hline
NFR-3.2 & Phân quyền theo vai trò (RBAC) &
  Hệ thống phân biệt ba vai trò: \textit{citizen, staff, admin}.
  API admin bị chặn hoàn toàn với người dùng không có quyền
  (trả về HTTP 403). \\
\hline
NFR-3.3 & Mã hóa mật khẩu &
  Mật khẩu người dùng phải được băm bằng thuật toán \textbf{bcrypt}
  trước khi lưu vào cơ sở dữ liệu. Không lưu mật khẩu dạng plaintext. \\
\hline
NFR-3.4 & Chính sách mật khẩu mạnh &
  Mật khẩu tối thiểu 8 ký tự, phải chứa chữ hoa, chữ số và ký tự đặc biệt.
  Áp dụng tại cả frontend (UX) và backend (validate). \\
\hline
NFR-3.5 & Bảo vệ file upload &
  Chỉ chấp nhận định dạng: PDF, TXT, JPG, JPEG, PNG, GIF.
  Giới hạn kích thước 16MB. Tên file được sanitize qua
  \texttt{secure\_filename()} trước khi lưu. \\
\hline
NFR-3.6 & Bảo vệ khỏi injection &
  Tất cả câu truy vấn SQL sử dụng parameterized query / SQLAlchemy ORM.
  Nghiêm cấm nối chuỗi trực tiếp vào câu SQL. \\
\hline
NFR-3.7 & Quản lý bí mật &
  Khóa API, chuỗi kết nối CSDL, JWT secret phải được lưu trong file
  \texttt{.env} (không commit vào git). Biến môi trường được đọc tại
  khởi động. \\
\hline
NFR-3.8 & Truyền thông mã hóa &
  Toàn bộ giao tiếp client–server trong môi trường production
  phải qua \textbf{HTTPS/TLS 1.2+}. Kết nối WebSocket qua \texttt{wss://}. \\
\hline
NFR-3.9 & Kiểm soát CORS &
  Backend cấu hình CORS chỉ cho phép origin từ domain frontend
  đã đăng ký; từ chối các origin không xác định. \\
\hline
NFR-3.10 & Bảo vệ dữ liệu nhạy cảm &
  Số CCCD, ngày sinh và thông tin cá nhân chỉ được trả về
  cho chính chủ tài khoản hoặc admin có thẩm quyền. \\
\hline
\end{tabularx}
\end{table}


% ─────────────────────────────────────────────────────────────
\subsection{Định nghĩa các tác nhân (Actors) trong hệ thống}
% ─────────────────────────────────────────────────────────────

Hệ thống E-Mapp có sáu tác nhân tham gia, được chia thành hai nhóm:
tác nhân con người (\textit{human actors}) và tác nhân hệ thống
(\textit{system actors}).

% ── Bảng tổng hợp ─────────────────────────────────────────────
\begin{table}[H]
\centering
\caption{Danh sách tác nhân trong hệ thống E-Mapp}
\label{tab:actors}
\begin{tabularx}{\textwidth}{|c|l|X|c|}
\hline
\textbf{Mã} & \textbf{Tác nhân} & \textbf{Mô tả} & \textbf{Loại} \\
\hline
A-01 & Công dân (Citizen)     & Người dùng cuối — sử dụng toàn bộ dịch vụ công     & Con người \\
A-02 & Nhân viên (Staff)      & Cán bộ cơ quan — điều hành quầy và hàng chờ        & Con người \\
A-03 & Quản trị viên (Admin)  & Quản lý toàn hệ thống, duyệt hồ sơ, cấu hình      & Con người \\
A-04 & AI Engine              & Mô-đun RAG + Voice AI xử lý ngôn ngữ tự nhiên      & Hệ thống \\
A-05 & VNeID                  & Hệ thống định danh điện tử quốc gia                & Hệ thống \\
A-06 & OpenStreetMap / OSRM   & Dịch vụ bản đồ và định tuyến mã nguồn mở          & Hệ thống \\
\hline
\end{tabularx}
\end{table}

% ── Chi tiết từng tác nhân ────────────────────────────────────
\subsubsection{A-01: Công dân (Citizen)}

\begin{description}
  \item[Định nghĩa:] Cư dân tỉnh Thanh Hóa hoặc bất kỳ cá nhân nào
    có nhu cầu sử dụng dịch vụ hành chính công qua nền tảng số.
  \item[Điều kiện tham gia:] Có số CCCD hợp lệ (12 chữ số); đăng ký
    và đăng nhập vào hệ thống. Một số chức năng cho phép khách
    (\textit{guest}) sử dụng mà không cần đăng nhập
    (xem bản đồ, lấy số hàng chờ, tra cứu thủ tục).
  \item[Phạm vi tương tác:] Đăng ký/đăng nhập, nộp và tra cứu hồ sơ,
    đặt lịch hẹn, lấy số thứ tự, sử dụng AI chatbot và voice,
    xem bản đồ, đánh giá dịch vụ.
  \item[Kênh truy cập:] Trình duyệt web trên máy tính và thiết bị di động.
\end{description}

\subsubsection{A-02: Nhân viên / Cán bộ (Staff)}

\begin{description}
  \item[Định nghĩa:] Cán bộ, nhân viên làm việc tại các cơ quan
    hành chính nhà nước sử dụng hệ thống để điều hành nghiệp vụ tại quầy.
  \item[Điều kiện tham gia:] Có tài khoản với vai trò \texttt{staff}
    được cấp bởi Admin.
  \item[Phạm vi tương tác:] Gọi số thứ tự tiếp theo, cập nhật trạng thái
    phục vụ, quản lý quầy giao dịch, xem danh sách vé hàng chờ
    và lịch hẹn của cơ quan.
  \item[Quyền hạn:] Cao hơn Công dân ở các chức năng điều hành hàng chờ
    và lịch hẹn; không có quyền xét duyệt hồ sơ hay quản trị hệ thống.
\end{description}

\subsubsection{A-03: Quản trị viên (Admin)}

\begin{description}
  \item[Định nghĩa:] Người quản lý và vận hành toàn bộ hệ thống,
    bao gồm cả việc xét duyệt hồ sơ và cấu hình hệ thống.
  \item[Điều kiện tham gia:] Có tài khoản với vai trò \texttt{admin}.
    Mọi route \texttt{/api/admin/*} đều kiểm tra vai trò tại middleware.
  \item[Phạm vi tương tác:] Duyệt/từ chối hồ sơ, quản lý tài khoản
    người dùng, CRUD địa điểm và danh mục thủ tục, xem thống kê
    \& analytics, cấu hình AI chatbot, quản lý hàng chờ.
  \item[Quyền hạn:] Cao nhất trong hệ thống; kế thừa toàn bộ
    quyền của Staff và Công dân.
\end{description}

\subsubsection{A-04: AI Engine (Tác nhân hệ thống)}

\begin{description}
  \item[Định nghĩa:] Mô-đun trí tuệ nhân tạo tích hợp trong backend,
    bao gồm hai thành phần chính:
    \begin{itemize}
      \item \textbf{RAG (Retrieval-Augmented Generation):} Sử dụng
        embedding tiếng Việt để tìm kiếm tài liệu và sinh câu trả lời
        theo ngữ cảnh qua \texttt{MultiRoleAgentGraph}.
      \item \textbf{Voice AI Backend:} Pipeline xử lý giọng nói tiếng Việt
        (STT → NLU/Gemini → TTS) hỗ trợ đặt lịch hẹn và hỏi đáp
        bằng giọng nói.
    \end{itemize}
  \item[Tương tác với:] Công dân (qua Chatbot và Voice),
    Google Cloud Speech API, Gemini API.
  \item[Vai trò:] Tác nhân thụ động — nhận yêu cầu từ người dùng,
    xử lý ngôn ngữ tự nhiên và trả về kết quả.
\end{description}

\subsubsection{A-05: VNeID (Tác nhân hệ thống)}

\begin{description}
  \item[Định nghĩa:] Hệ thống định danh điện tử quốc gia của Việt Nam
    do Cục Cảnh sát Quản lý Hành chính về Trật tự Xã hội (C06)
    vận hành.
  \item[Vai trò trong hệ thống:] Xác thực tính hợp lệ của số CCCD
    trong quá trình đăng ký tài khoản. Khi người dùng chọn
    đăng ký qua VNeID, hệ thống gọi API \texttt{verify\_vneid()}
    để kiểm tra và lấy thông tin định danh.
  \item[Tương tác với:] Chức năng đăng ký tài khoản (FR-1.2).
\end{description}

\subsubsection{A-06: OpenStreetMap / OSRM (Tác nhân hệ thống)}

\begin{description}
  \item[Định nghĩa:] Tập hợp các dịch vụ bản đồ mã nguồn mở miễn phí:
    \begin{itemize}
      \item \textbf{OpenStreetMap + Nominatim:} Dữ liệu bản đồ và
        dịch vụ geocoding/reverse-geocoding/autocomplete.
      \item \textbf{OSRM (Open Source Routing Machine):} Dịch vụ
        tính toán lộ trình đa phương thức (đi bộ, xe máy, ô tô).
    \end{itemize}
  \item[Lý do chọn:] Hoàn toàn miễn phí, không yêu cầu API key,
    phù hợp triển khai trong dự án chính phủ không có chi phí
    bản quyền bản đồ.
  \item[Tương tác với:] Module bản đồ (FR-5.1 đến FR-5.6).
\end{description}
